\setcounterpageref{ex}{0}
\TaoPhieuKiemTra{16}

\setcounter{ex}{0}
\setcounter{bt}{0}

\textbf{\large I. PHẦN TRẮC NGHIỆM (3,0 điểm)}

{\itshape Khoanh tròn vào chữ cái đặt trước câu trả lời đúng. Mỗi câu đúng được 0,5 điểm.}

\begin{ex} % Câu 1
Trong số $95$, chữ số $5$ thuộc hàng
\choice
{Hàng chục}
{Hàng trăm}
{Hàng chục và đơn vị}
{\True Hàng đơn vị}
\loigiai{
Trong số $95$, chữ số $5$ thuộc hàng đơn vị, chữ số $9$ thuộc hàng chục.\\
\textbf{Đáp án đúng: D}
}
\end{ex}

\begin{ex} % Câu 2
Trong phép tính $80 - 30 = 50$, số $30$ được gọi là
\choice
{Số bị trừ}
{Hiệu}
{\True Số trừ}
{Số hạng}
\loigiai{
Số $30$ đứng sau dấu trừ nên gọi là Số trừ.\\
\textbf{Đáp án đúng: C}
}
\end{ex}

\begin{ex} % Câu 3
Gang tay của mẹ dài khoảng
\choice
{$2\text{ kg}$}
{\True $2\text{ dm}$}
{$2\text{ }l$}
{$2\text{ cm}$}
\loigiai{
Gang tay dài khoảng $2\text{ dm}$ (tương đương $20\text{ cm}$).\\
\textbf{Đáp án đúng: B}
}
\end{ex}

\begin{ex} % Câu 4
Đoạn thẳng $OA$ dài $6\text{ cm}$, đoạn thẳng $OB$ dài $9\text{ cm}$. Đoạn thẳng $OA$ ngắn hơn đoạn $OB$ là
\choice
{\True $3\text{ cm}$}
{$2\text{ cm}$}
{$4\text{ cm}$}
{$15\text{ cm}$}
\loigiai{
$9 - 6 = 3\text{ cm}$.\\
\textbf{Đáp án đúng: A}
}
\end{ex}

\begin{ex} % Câu 5
Điền dấu: $67 - 24 \dots 45$.
\choice
{$>$}
{$=$}
{\True $<$}
{Không so sánh được}
\loigiai{
$67 - 24 = 43$. Vì $43 < 45$ nên điền dấu $<$.\\
\textbf{Đáp án đúng: C}
}
\end{ex}

\begin{ex} % Câu 6
Tính: $13 - 6 + 5 = \dots$
\choice
{$11$}
{$13$}
{$10$}
{\True $12$}
\loigiai{
$13 - 6 = 7$; $7 + 5 = 12$.\\
\textbf{Đáp án đúng: D}
}
\end{ex}

\vspace{4mm}
\textbf{\large II. PHẦN TỰ LUẬN (7,0 điểm)}

{\itshape Thí sinh trình bày chi tiết bài làm vào giấy kiểm tra.}

\begin{bt}[1,0 điểm]
Tính nhẩm:
\begin{multicols}{2}
\begin{enumerate}[a)]
  \item $8 + 5$
  \item $7 + 6$
  \item $13 - 5$
  \item $15 - 6$
\end{enumerate}
\end{multicols}
\loigiai{
\begin{itemize}
  \item $8 + 5 = 13$
  \item $7 + 6 = 13$
  \item $13 - 5 = 8$
  \item $15 - 6 = 9$
\end{itemize}
}
\end{bt}

\begin{bt}[1,5 điểm]
Đặt tính rồi tính:
\begin{multicols}{3}
\begin{enumerate}[a)]
  \item $25 + 55$
  \item $87 - 29$
  \item $48 + 44$
\end{enumerate}
\end{multicols}
\loigiai{
\begin{itemize}
  \item $25 + 55 = 80$
  \item $87 - 29 = 58$
  \item $48 + 44 = 92$
\end{itemize}
}
\end{bt}

\begin{bt}[1,0 điểm]
Tính giá trị biểu thức kèm đơn vị đo:
\begin{enumerate}[a)]
  \item $32\text{ kg} + 7\text{ kg} - 10\text{ kg}$
  \item $41\text{ }l - 12\text{ }l + 6\text{ }l$
\end{enumerate}
\loigiai{
\begin{enumerate}[a)]
  \item $32\text{ kg} + 7\text{ kg} - 10\text{ kg} = 39\text{ kg} - 10\text{ kg} = 29\text{ kg}$.
  \item $41\text{ }l - 12\text{ }l + 6\text{ }l = 29\text{ }l + 6\text{ }l = 35\text{ }l$.
\end{enumerate}
}
\end{bt}

\begin{bt}[1,5 điểm]
Bạn Khánh có $31$ quyển vở. Bạn Linh có ít hơn bạn Khánh $8$ quyển vở. Hỏi bạn Linh có bao nhiêu quyển vở?
\loigiai{
Số quyển vở bạn Linh có là:
\[ 31 - 8 = 23 \text{ (quyển vở)} \]
\textbf{Đáp số:} $23$ quyển vở.
}
\end{bt}

\begin{bt}[1,0 điểm]
Một thùng chứa đầy $51\text{ }l$ dầu ăn. Buổi sáng cửa hàng đã bán ra $12\text{ }l$. Hỏi trong thùng còn lại bao nhiêu lít dầu ăn?
\loigiai{
Trong thùng còn lại số lít dầu ăn là:
\[ 51 - 12 = 39\text{ }l \]
\textbf{Đáp số:} $39\text{ }l$ dầu ăn.
}
\end{bt}

\begin{bt}[1,0 điểm]
Điền số thích hợp vào chỗ chấm: $2, 5, 8, 11, \dots, \dots$
\loigiai{
Quy luật: Dãy số tăng đều $3$ đơn vị.\\
Số tiếp theo: $11 + 3 = 14$.\\
Số tiếp theo nữa: $14 + 3 = 17$.\\
\textbf{Đáp số:} $14$ và $17$.
}
\end{bt}
